\documentclass[tikz,border=10pt]{standalone}
\usepackage[UTF8]{ctex}
\usetikzlibrary{arrows.meta,backgrounds,fit,positioning,shadows.blur,shapes.geometric}

\definecolor{ink}{HTML}{1F2937}
\definecolor{muted}{HTML}{5B6472}
\definecolor{lineblue}{HTML}{2563EB}
\definecolor{frontendfill}{HTML}{EFF6FF}
\definecolor{backendfill}{HTML}{F0FDF4}
\definecolor{datafill}{HTML}{FFF7ED}
\definecolor{actorfill}{HTML}{F8FAFC}
\definecolor{groupborder}{HTML}{CBD5E1}

\tikzset{
  font=\sffamily,
  >={Latex[length=3mm,width=2mm]},
  actor/.style={draw=ink, rounded corners=2pt, fill=actorfill,
    minimum width=28mm, minimum height=13mm, align=center,
    line width=.8pt, font=\bfseries\sffamily},
  component/.style={draw=ink, rounded corners=2pt, fill=white,
    minimum width=40mm, minimum height=18mm, text width=36mm,
    align=center, inner sep=4pt, line width=.8pt,
    font=\small\sffamily, blur shadow={shadow opacity=.10}},
  datastore/.style={cylinder, shape border rotate=90, aspect=.27,
    draw=ink, fill=white, minimum width=40mm, minimum height=19mm,
    align=center, line width=.8pt, font=\small\sffamily},
  dependency/.style={-Latex, draw=muted, line width=.8pt},
  request/.style={-Latex, draw=lineblue, line width=1.1pt},
  labeltext/.style={font=\scriptsize\sffamily, fill=white,
    inner sep=1.5pt, text=muted},
  group/.style={draw=groupborder, rounded corners=5pt,
    line width=.8pt, inner sep=2mm},
  grouptitle/.style={font=\bfseries\small\sffamily, text=ink,
    fill=white, inner sep=2pt},
  note/.style={draw=groupborder, rounded corners=2pt, fill=actorfill,
    text width=325mm, align=left, inner sep=5pt,
    font=\small\sffamily, text=ink}
}

\newcommand{\componentglyph}[1]{%
  \begin{scope}[shift={(#1.north east)},xshift=-2.5mm,yshift=-1.5mm,
      draw=ink,line width=.55pt]
    \draw (-3mm,-1.7mm) rectangle (0,1.7mm);
    \draw (-4.2mm,0.6mm) rectangle (-2.5mm,1.55mm);
    \draw (-4.2mm,-1.55mm) rectangle (-2.5mm,-0.6mm);
  \end{scope}%
}

\newcommand{\diagramtitle}[2]{%
  \node[font=\bfseries\LARGE\sffamily, text=ink] (title) at (15.5,9.5) {#1};
  \node[below=1.5mm of title, font=\small\sffamily, text=muted] {#2};
}

\newcommand{\groupcaption}[2]{%
  \node[grouptitle, anchor=south west]
    at ([xshift=2mm,yshift=1mm]#1.north west) {#2};
}

\begin{document}
\begin{tikzpicture}
  \diagramtitle{UC04 评论与赞踩——UML 组件图}
    {展示评论发布、视频／评论赞踩、状态切换及互动通知所依赖的组件}

  \node[actor] (user) at (0,5.7) {登录用户};
  \node[actor] (author) at (0,1.8) {内容作者\\（被通知方）};
  \node[component] (detail) at (4,5.7) {\textbf{VideoDetailView}\\\footnotesize 视频互动入口};
  \node[component] (commentItem) at (4,1.8) {\textbf{CommentItem}\\\footnotesize 评论列表与互动状态};
  \node[component] (client) at (8.5,3.75) {\textbf{comment.js + reaction.js}\\\footnotesize 评论与赞踩请求};

  \node[component] (security) at (13,3.75) {\textbf{SecurityConfig + JWT Filter}\\\footnotesize 登录身份校验};
  \node[component] (commentCtl) at (17.5,5.7) {\textbf{CommentController}\\\footnotesize 评论新增与查询};
  \node[component] (reactionCtl) at (17.5,1.8) {\textbf{ReactionController}\\\footnotesize 视频／评论赞踩};
  \node[component] (commentSvc) at (22,5.7) {\textbf{CommentService}\\\footnotesize 内容校验与评论持久化};
  \node[component] (reactionSvc) at (22,1.8) {\textbf{ReactionService}\\\footnotesize 切换、取消与计数};

  \node[component] (mapper) at (26.5,4.6) {\textbf{Mapper 层}\\\footnotesize CommentMapper\\Reaction Mappers};
  \node[component] (notify) at (22,-0.6) {\textbf{NotificationService}\\\footnotesize 互动事件通知};
  \node[datastore] (mysql) at (31,3.75) {\textbf{MySQL 8}\\\footnotesize comments / reactions\\notifications};

  \draw[request] (user) -- node[labeltext,above]{评论／赞踩} (detail);
  \draw[request] (author) -- node[labeltext,above]{查看结果} (commentItem);
  \draw[request] (detail) -- (client);
  \draw[request] (commentItem) -- (client);
  \draw[request] (client) -- node[labeltext,above]{HTTP + JWT} (security);
  \draw[request] (security) to[out=30,in=180] (commentCtl);
  \draw[request] (security) to[out=-30,in=180] (reactionCtl);
  \draw[dependency] (commentCtl) -- (commentSvc);
  \draw[dependency] (reactionCtl) -- (reactionSvc);
  \draw[dependency] (commentSvc) to[out=0,in=150] (mapper);
  \draw[dependency] (reactionSvc) to[out=0,in=-150] (mapper);
  \draw[dependency] (reactionSvc) -- (notify);
  \draw[dependency] (mapper) -- (mysql);
  \draw[dependency] (notify.east) to[out=20,in=-90] (mapper.south);

  \foreach \n in {detail,commentItem,client,security,commentCtl,reactionCtl,commentSvc,reactionSvc,mapper,notify}{\componentglyph{\n}}
  \begin{scope}[on background layer]
    \node[group, fill=frontendfill, fit=(detail)(commentItem)(client)] (fg) {};
    \node[group, fill=backendfill, fit=(security)(commentCtl)(reactionCtl)(commentSvc)(reactionSvc)(notify)] (bg) {};
    \node[group, fill=datafill, fit=(mapper)(mysql)] (dg) {};
  \end{scope}
  \groupcaption{fg}{Vue 3 前端}
  \groupcaption{bg}{Spring Boot 后端}
  \groupcaption{dg}{数据与通知基础设施}
  \node[note, anchor=north west] at (-1.4,-2.0) {\textbf{核心规则：}\quad 同一用户对同一视频或评论只保留一种反应；重复点击可取消，点击相反反应则切换；空评论和不存在的目标必须拒绝。};
\end{tikzpicture}
\end{document}
